\documentclass[12pt,a4paper]{article}
\usepackage{ctex}
\usepackage{geometry}
\usepackage{graphicx}
\usepackage{amsmath}
\usepackage{booktabs}
\usepackage{multirow}
\usepackage{caption}
\usepackage{subcaption}
\usepackage{float}
\usepackage{tikz}
\usepackage{pgfplots}
\pgfplotsset{compat=1.16}

\geometry{left=2.5cm,right=2.5cm,top=2.5cm,bottom=2.5cm}

\title{\textbf{磁滞回线的测量实验报告}}
\author{}
\date{}

\newcommand{\recordblank}[1]{\par\vspace*{#1}\par}

\usepackage{xcolor}
\definecolor{questionblue}{RGB}{39,82,201}
\usepackage[most]{tcolorbox}
\newtcolorbox{questionbox}{
  top=2pt,bottom=2pt,left=5pt,right=5pt,
  arc=1mm,boxrule=0.5pt,
  colframe=questionblue,
  colback=questionblue!3!white,
  breakable
}

\begin{document}

\begin{center}
    {\LARGE\textbf{磁滞回线的测量}}\\[0.5cm]
    {\Large 实验报告}\\[1cm]
\end{center}

\begin{center}
    \textbf{姓名：}\underline{\makebox[4cm][c]{}} \quad
    \textbf{温度：}\underline{\makebox[2.2cm][c]{}} \quad
    \textbf{湿度：}\underline{\makebox[2cm][c]{}}
\end{center}

\begin{center}
    \textbf{实验日期：}\underline{\makebox[3.5cm][c]{}} \quad
    \textbf{星期：}\underline{\makebox[2cm][c]{}} \quad
    \textbf{节次：}\underline{\makebox[3.5cm][c]{}} \\[0.3cm]
    \textbf{座位号：}\underline{\makebox[3cm][c]{}}
\end{center}

\section*{实验仪器}
\begin{enumerate}
    \item 电源 1台
    \item 功率放大电路板 1件
    \item 环形变压器 1个（硅钢材料）
    \item myDAQ数据采集单元 1个
    \item 实验电路板 1件
    \item 导线若干
\end{enumerate}

\section*{实验内容}

\subsection*{1. 测试环形变压器在三组不同频率和幅值激励信号下的饱和磁滞回线}

\subsubsection*{实验记录}
\recordblank{5.9cm}
\subsection*{2.实验结果}
\recordblank{18cm}
\clearpage
\recordblank{8.4cm}
\subsection*{2. 三种材料的磁滞回线对比测试}

\subsubsection*{实验记录}
\recordblank{5.0cm}
\subsubsection*{实验结果}
\recordblank{9.8cm}
\section*{实验结论}
\subsection*{1. 实验总结}
\subsubsection*{(1) 激励信号参数对磁滞回线的影响}
\recordblank{3.5cm}
\subsubsection*{(2) 三种材料的磁特性对比}
\recordblank{3.5cm}
\subsection*{2. 实验过程遇到的问题}
\recordblank{5.0cm}
\subsection*{3. 预习思考题}

\begin{questionbox}
思考题1：对于铁磁性物质，为什么缠上线圈，通上电流，就能在缠绕物质中产生非常强的磁感应强度？微观物理机制是什么？
\end{questionbox}
\recordblank{4cm}
\begin{questionbox}
思考题2：如何选择初级线圈的匝数、次级线圈匝数、电流采样电阻R1？如何设置交流电压频率和电压幅值，两者间有何关系，如何避免线圈中电流过大？
\end{questionbox}
\recordblank{4cm}
\subsection*{4. 实验后思考题}

\begin{questionbox}
思考题1：各向异性磁电阻(AMR)传感器为何要退磁？
\end{questionbox}
\recordblank{4cm}
\begin{questionbox}
思考题2：为什么输入、输出正弦波会发生畸变？什么情况下畸变比较小？
\end{questionbox}
\recordblank{4cm}
\begin{questionbox}
思考题3：铁磁材料的磁导率在什么条件下才可以近似为一常数？
\end{questionbox}
\recordblank{4cm}
\noindent\textbf{成绩：}\underline{\hspace{3cm}} \hspace{1.5cm} \textbf{教师签名：}\underline{\hspace{3cm}} \hspace{1.5cm} \textbf{日期：}\underline{\hspace{3cm}}

\section*{附录：实验记录图片}


\recordblank{7cm}


\recordblank{7cm}


\recordblank{7cm}


\recordblank{7cm}

\end{document}
